% Part: first-order-logic
% Chapter: lindstrom
% Section: ls-property

\documentclass[../../../include/open-logic-section]{subfiles}

\begin{document}

\olfileid{mod}{lin}{lsp}

\olsection{কম্প্যাক্টনেস ও লোয়েনহাইম--স্কোলেম ধর্ম}

এখন বিমূর্ত যুক্তিবিদ্যার ক্ষেত্রে কম্প্যাক্টনেস ও
লোয়েনহাইম--স্কোলেম ধর্মের স্বাভাবিক সম্প্রসারণগুলি দিচ্ছি।

\begin{defn}
একটি বিমূর্ত যুক্তিবিদ্যা $\tuple{L, \models_L}$-এর
\emph{কম্প্যাক্টনেস ধর্ম} আছে যদি $L(\Lang{L})$-এর !!{sentence}s-এর
প্রতিটি সেট $\Gamma$ সন্তোষণীয় হয় যখনই প্রতিটি সসীম
 $\Gamma_0 \subseteq \Gamma$ সন্তোষণীয় হয়।
\end{defn}

\begin{defn}
$\tuple{L, \models_L}$-এর \emph{নিম্নমুখী লোয়েনহাইম--স্কোলেম
ধর্ম} আছে যদি যেকোনো সন্তোষণীয় $\Gamma$-র একটি !!a{enumerable}
মডেল থাকে।
\end{defn}

\olref[bas][pis]{defn:partialisom}-এ সংজ্ঞায়িত আংশিক সমরূপতার ধারণাটি
সম্পূর্ণ “বীজগাণিতিক” (অর্থাৎ ভাষার !!{sentence}s-এর কোনো উল্লেখ ছাড়াই,
শুধু !!{language}~$\Lang{L}$-এর !!{structure}s-এ দেওয়া ধ্রুবকগুলির
সাপেক্ষে সংজ্ঞায়িত), তাই এটি বিমূর্ত যুক্তিবিদ্যার ক্ষেত্রেও প্রযোজ্য।
প্রথম-ক্রমের যুক্তিবিদ্যার ক্ষেত্রে আমরা
\olref[bas][pis]{thm:p-isom2} থেকে জানি যে দুটি !!{structure} (দুটি !!{structure}s) আংশিকভাবে
সমরূপ হলে তারা মৌলিকভাবে সমতুল্য। এই প্রমাণটি বিমূর্ত যুক্তিবিদ্যায়
চলে না, কারণ যেকোনো $!E \in L(\Lang{L})$-এর জন্য !!{formula}s-এর ওপর
আবেশ প্রযোজ্য নাও হতে পারে; তবু লোয়েনহাইম--স্কোলেম ধর্ম সত্য হলে
উপপাদ্যটি সত্য।

\begin{thm}
\ollabel{thm:abstract-p-isom}
ধরা যাক $\tuple{L, \models_L}$ একটি স্বাভাবিক যুক্তিবিদ্যা যার
লোয়েনহাইম--স্কোলেম ধর্ম আছে। তাহলে যেকোনো দুটি !!{structure} যারা
আংশিকভাবে সমরূপ, তারা $\tuple{L,\models_L}$-এ মৌলিকভাবে সমতুল্য।
\end{thm}

\begin{proof}
ধরি $\Struct{M} \simeq_p \Struct{N}$, কিন্তু কোনো একটি $!E$-এর জন্য
$\Struct{M} \models_L !E$ এবং $\Struct{N} \not\models_L !E$।
সমরূপতা ধর্ম অনুসারে ধরে নিতে পারি যে $\Domain{M}$ ও $\Domain{N}$
পরস্পর বিচ্ছিন্ন, এবং প্রসারণ ধর্ম অনুসারে ধরে নিতে পারি যে সসীম
!!{language}~$\Lang{L}$-এর জন্য $!E \in L(\Lang{L})$।
$\Struct{M}$ ও $\Struct{N}$-এর মধ্যকার আংশিক সমরূপগুলির একটি সেট
$\mathcal{I}$ নিই; সাধারণতা না হারিয়ে আরও ধরে নিই, $p \in \mathcal{I}$
এবং $q \subseteq p$ হলে $q \in \mathcal{I}$।

$\Domain{M}^{<\omega}$ হলো $\Domain{M}$-এর !!{element}s-উপাদানগুলির সসীম
ক্রমগুলির সেট। $\Domain{M}^{<\omega}$-এর ওপর $S$-কে ত্রিস্থানিক
সম্পর্ক ধরি, যা সংযোজন নির্দেশ করে: $\mathbf{a}, \mathbf{b},
\mathbf{c} \in \Domain{M}^{<\omega}$ হলে $S(\mathbf{a}, \mathbf{b},
\mathbf{c})$ সত্য ঠিক তখনই যখন $\mathbf{c}$ হলো $\mathbf{a}$ ও
$\mathbf{b}$-এর সংযোজন; এবং $T$-কে এমন ত্রিস্থানিক সম্পর্ক ধরি যে
$b \in M$ এবং $\mathbf{a}, \mathbf{c} \in \Domain{M}^{<\omega}$-এর
জন্য $T(\mathbf{a}, b, \mathbf{c})$ সত্য ঠিক তখনই যখন
$\mathbf{a} = a_1, \dots a_n$ এবং $\mathbf{c} = a_1, \dots a_n, b$।
নতুন 3-place !!{predicate}s $P$ ও $Q$ নিই এবং
$\Struct{M}^*$ কাঠামো গঠন করি, যার universe
$\Domain{M} \cup \Domain{M}^{<\omega}$, যাতে $\Struct{M}$ একটি
উপকাঠামো এবং $P$, $Q$-এর ব্যাখ্যা যথাক্রমে $S$, $T$ সংযোজন-সম্পর্ক
(তাই $\Struct{M}^*$ হলো !!{language}
$\Lang{L} \cup \{ P, Q\}$-এর একটি কাঠামো)।

$\Domain{N}^{<\omega}$, $S'$, $T'$, $P'$, $Q'$ এবং
$\Struct{N}^*$ একইভাবে সংজ্ঞায়িত করি। যেহেতু অনুমান অনুযায়ী
$\Struct{M} \simeq_p \Struct{N}$, তাই $\Domain{M}^{<\omega}$ ও
$\Domain{N}^{<\omega}$-এর মধ্যে এমন একটি সম্পর্ক $I$ আছে যাতে
$I(\mathbf{a}, \mathbf{b})$ সত্য ঠিক তখনই যখন $\mathbf{a}$ ও
$\mathbf{b}$ সমরূপ এবং \olref[bas][pis]{defn:partialisom}-এর
আগ-পিছু শর্ত পূরণ করে। এই $\Struct{M}$-কে এখন এমন কাঠামো ধরি যার
!!{domain} হলো $\Struct{M}^*$ ও $\Struct{N}^*$-এর !!{domain}s-এর
সংযুক্তি, যাতে $\Struct{M}^*$ ও $\Struct{N}^*$ উপ!!{structure}s, এবং
!!{language}-এ অতিরিক্ত একটি দ্বিস্থানিক !!{predicate}~$R$ আছে যার
ব্যাখ্যা $I$; পাশাপাশি $\Domain{M}^*$ ও $\Domain{N}*$ নির্দেশকারী
!!{domain}s নির্দেশকারী !!{predicate}s-ও আছে।

\begin{figure}[h]
  \centering
  \begin{tikzpicture}[node distance=2cm, auto, thick, >=stealth']
    \draw [rounded corners] (0,0) -- (8,0) -- (8,4) -- (0,4) --  cycle;
    \draw (2,2) circle (0.5cm);
    \draw (2,2) circle (1.25cm);
    \draw (6,2) circle (0.5cm);
    \draw (6,2) circle (1.25cm);
    \path node at (0.75,3.5) {\large $\Struct{M}$};
    \path node at (2,2) {\large $\Struct{M}$};
    \path node at (6,2) {\large $\Struct{N}$};
    \path node at (3.5,1) {\large $\Struct{M}^*$};
    \path node at (7.5,1) {\large $\Struct{N}^*$};
    \node (Idom) at (2.8,2) {};
    \node (Irng) at (5.2,2) {};
    \draw[<->, bend left] (Idom) to node {\large $I$} (Irng) ;
  \end{tikzpicture}
  \caption{অভ্যন্তরীণ আংশিক সমরূপতা-সহ !!{structure}~$\Struct{M}$।}
\end{figure}

গুরুত্বপূর্ণ পর্যবেক্ষণটি হলো, !!{structure}~$\Struct{M}$-এর
!!{language}-এ একটি \emph{প্রথম-ক্রমের} !!{sentence} $!D_1$ আছে, যা
$\Struct{M}$-এ সত্য এবং বলে যে $\Struct{M} \models_L !E$ ও
$\Struct{N} \not\models_L !E$ (এর জন্য Relativization Property দরকার),
এবং এমন আরেকটি \emph{প্রথম-ক্রমের} !!{sentence} $!D_2$ আছে, যা
$\Struct{M}$-এ সত্য এবং বলে যে আংশিক সমরূপতা~$I$-এর মাধ্যমে
$\Struct{M} \simeq_p \Struct{N}$। লোয়েনহাইম--স্কোলেম ধর্ম অনুসারে,
$!D_1$ ও $!D_2$ যৌথভাবে সত্য এমন একটি !!a{enumerable} মডেল
$\Struct{M}_0$ পাওয়া যায়, যাতে আংশিকভাবে সমরূপ উপ!!{structure}s
$\Struct{M}_0$ ও $\Struct{N}_0$ আছে এবং
$\Struct{M}_0 \models_L !E$, কিন্তু $\Struct{N}_0
\not\models_L !E$। অথচ !!{enumerable} আংশিকভাবে সমরূপ !!{structure}s
আসলে সমরূপ, \olref[bas][pis]{thm:p-isom1} অনুসারে; এটি স্বাভাবিক
বিমূর্ত যুক্তিবিদ্যার সমরূপতা ধর্মের বিরোধী।
\end{proof}

\end{document}
